\section{Endpoint recovery and the fixed-shift strength threshold}
\label{sec:endpoint-strength}

The ratio transform does not select $m=1$ by merely taking a long sharp
frequency interval.  A fixed scale window does select it, but that window
uses a coherent superposition of all frequencies.

Assume in this section that
\begin{equation}
 \supp W\subset[\alpha,\beta],
 \qquad \beta<4\alpha,
 \qquad D<\alpha X.
 \label{eq:endpoint-support-hypothesis}
\end{equation}
Choose $V\in C_c^\infty(\R)$ such that
\begin{equation}
 V(x)=1\quad(\log\alpha\leq x\leq\log\beta),
 \qquad
 V(x)=0\quad(x\leq\log(\beta/4)).
 \label{eq:endpoint-window}
\end{equation}

\begin{proposition}[Exact one-Mellin endpoint selector]
\label{prop:one-Mellin-endpoint-selector}
Under \eqref{eq:endpoint-support-hypothesis},
\begin{align}
 \sum_n u(n)\Lambda(n+h)W(n/X)
 &=\frac1{2\pi}\int_\R
 \widehat V(t)X^{-it}\cT_{h,D}(X;t)dt,
 \label{eq:raw-endpoint-selector}\\
 \sum_n u(n)(\Lambda(n+h)-1)W(n/X)
 &=\frac1{2\pi}\int_\R
 \widehat V(t)X^{-it}\cR_{h,D}(X;t)dt,
 \label{eq:centered-endpoint-selector}
\end{align}
where $\widehat V(t)=\int_\R V(x)e^{-itx}dx$.
\end{proposition}

\begin{proof}
If $r=mn$ lies in the support of $W(r/X)$, then
\begin{equation}
 \log(n/m)-\log X=\log(r/X)-2\log m.
 \label{eq:endpoint-coordinate}
\end{equation}
For $m=1$ this lies in $[\log\alpha,\log\beta]$, where $V=1$.  For
$m\geq2$ it is at most $\log\beta-2\log2=\log(\beta/4)$, where $V=0$.
Thus inserting $V(\log(n/m)-\log X)$ selects precisely $m=1$.
The cutoff $D<\alpha X$ removes none of those rows.  Fourier inversion of
$V$ gives \eqref{eq:raw-endpoint-selector}; subtracting the same identity
for \eqref{eq:intro-source-model} gives
\eqref{eq:centered-endpoint-selector}.
\end{proof}

\begin{remark}[General product supports]
Every compact subinterval of $(0,\infty)$ is a finite union of intervals
of multiplicative width less than four.  A smooth partition of unity in
the product variable therefore reduces a general $W$ to finitely many
selectors of the form above.
\end{remark}

\subsection{Truncating the coherent selector}

The total variation of either factor-scale measure in
\eqref{eq:raw-endpoint-selector}--\eqref{eq:centered-endpoint-selector}
satisfies
\begin{equation}
 \sum_{\substack{n>D\\m\geq1}}
 |u(n)|\bigl(1+\Lambda(mn+h)\bigr)|W(mn/X)|
 \ll_{h,W}X\log(2X).
 \label{eq:endpoint-total-variation}
\end{equation}
One way to see this is to put $r=mn$, use
\[
 \sum_{n\mid r}|u(n)|\leq\log r+\tau(r),
\]
and apply the prime number theorem to the logarithmic part and the
standard shifted Titchmarsh upper bound to the divisor part
\citep{AssingBlomerLi2021}.

Since $V\in C_c^\infty$, for every $A>0$,
\begin{equation}
 \left|\int_{|t|>T}\widehat V(t)X^{-it}
 \cR_{h,D}(X;t)dt\right|
 \ll_{A,h,V,W}X\log(2X)T^{-A}.
 \label{eq:endpoint-Fourier-tail}
\end{equation}
The same estimate holds for the raw transform and for the source model.
Thus any positive power of $\log X$ is a sufficient truncation length for
an $o(X)$ approximation, after choosing $A$ large.  The hard part is not
the smooth Fourier tail; it is the accuracy of the arithmetic model on
the retained band.

\subsection{An exact \texorpdfstring{$L^2$}{L2} strength test}

For a function $F$ for which the integral converges, define
\begin{equation}
 \mathscr L_X(F):=\frac1{2\pi}\int_\R
 \widehat V(t)X^{-it}F(t)dt.
 \label{eq:endpoint-functional}
\end{equation}

\begin{proposition}[Endpoint discrepancy forces band energy]
\label{prop:endpoint-strength-test}
Let $S$ and $M$ be two transforms and let
\begin{equation}
 \operatorname{Tail}_T(S-M)
 :=\frac1{2\pi}
 \left|\int_{|t|>T}\widehat V(t)X^{-it}(S(t)-M(t))dt\right|.
 \label{eq:model-tail}
\end{equation}
Then
\begin{equation}
 \|S-M\|_{L^2[-T,T]}
 \geq\frac{2\pi}{\|\widehat V\|_{L^2(\R)}}
 \left(
 |\mathscr L_X(S)-\mathscr L_X(M)|
 -\operatorname{Tail}_T(S-M)
 \right)_+.
 \label{eq:endpoint-strength-lower-bound}
\end{equation}
Consequently, if the endpoint discrepancy is $\gg X$ and the Fourier tail
is $o(X)$, then
\begin{equation}
 \int_{-T}^{T}|S(t)-M(t)|^2dt\gg X^2.
 \label{eq:endpoint-X2-floor}
\end{equation}
Conversely, an absolute band error $o(X^2)$ together with an $o(X)$ tail
forces the two endpoint functionals to agree up to $o(X)$.
\end{proposition}

\begin{proof}
Split \eqref{eq:endpoint-functional} at $|t|=T$, apply the triangle
inequality, and then apply Cauchy--Schwarz on $[-T,T]$:
\[
 |\mathscr L_X(S)-\mathscr L_X(M)|
 \leq\frac{\|\widehat V\|_2}{2\pi}
 \|S-M\|_{L^2[-T,T]}
 +\operatorname{Tail}_T(S-M).
\]
Rearrangement proves \eqref{eq:endpoint-strength-lower-bound}.  The two
consequences are immediate.
\end{proof}

The normalization is important.  A statement
\[
 \int_{-T}^{T}|S-M|^2dt=o(TX^2)
\]
is only a root-mean-square statement.  When $T\to\infty$, Cauchy--Schwarz
allows an $o(X\sqrt T)$ coherent error and does not recover the endpoint.
The absolute $o(X^2)$ threshold in
\cref{prop:endpoint-strength-test} is strictly stronger.

\subsection{The Hardy--Littlewood boundary}

For a nonzero integer shift $h$, write the usual prime-pair singular
series as
\begin{equation}
 \mathfrak S(h):=
 \begin{cases}
  2C_2\displaystyle\prod_{\substack{p\mid h\\p>2}}
  \dfrac{p-1}{p-2},&2\mid h,\\[0.8em]
  0,&2\nmid h,
 \end{cases}
 \qquad
 C_2:=\prod_{p>2}\left(1-\frac1{(p-1)^2}\right).
 \label{eq:prime-pair-singular-series}
\end{equation}
This is the constant in the corresponding Hardy--Littlewood prediction
\citep{HardyLittlewood1923}; no such prediction is assumed below.

For the raw prime-target transform, the selected endpoint is
\begin{equation}
 \begin{aligned}
 \mathscr L_X(\cT_{h,D})
 &=\sum_n(\Lambda(n)-1)\Lambda(n+h)W(n/X)\\
 &=\sum_n\Lambda(n)\Lambda(n+h)W(n/X)
 -X I_0(W)+o_{h,W}(X),
 \end{aligned}
 \label{eq:endpoint-HL-reduction}
\end{equation}
by the shifted prime number theorem.  Hence, for one fixed admissible even
shift and one fixed $W$, evaluating \eqref{eq:endpoint-HL-reduction} with
the predicted main term is equivalent to the corresponding single-weight
Hardy--Littlewood assertion.  The centered endpoint has the same predicted
leading term because
$\sum_n(\Lambda(n)-1)W(n/X)=o_W(X)$.

It follows from \cref{prop:endpoint-strength-test} that a fixed-shift band
model with absolute $o(X^2)$ error, a compatible filtered model value, and
$o(X)$ tails would already contain that Hardy--Littlewood endpoint.  The
implication is one-way: the endpoint controls one filtered functional, not
the entire band.

There is an unconditional parity check on the natural source model.

\begin{corollary}[Odd-shift model floor]
\label{cor:odd-shift-model-floor}
Assume $I_0(W)\ne0$ and fix an odd integer $h$.  Let
$S=\cT_{h,D}$ and $M=\cM_D$.  If $T\to\infty$ so that the tail in
\eqref{eq:model-tail} is $o(X)$---for example any fixed positive power of
$\log X$---then
\begin{equation}
 \int_{-T}^{T}|\cT_{h,D}(X;t)-\cM_D(X;t)|^2dt
 \gg_{h,V,W}X^2.
 \label{eq:odd-shift-band-floor}
\end{equation}
\end{corollary}

\begin{proof}
If $h$ is odd, two prime-power-supported values $n$ and $n+h$ can both
occur only when one is a power of two.  There are $O(\log X)$ such
possibilities, so
\[
 \sum_n\Lambda(n)\Lambda(n+h)W(n/X)=o_{h,W}(X).
\]
The prime number theorem applied to the two linear terms in the centered
endpoint then gives
\[
 \mathscr L_X(\cT_{h,D}-\cM_D)
 =-X I_0(W)+o_{h,W}(X).
\]
Apply \cref{prop:endpoint-strength-test} and
\eqref{eq:endpoint-Fourier-tail}.
\end{proof}

\begin{remark}[Why this is not a twin-prime lower bound]
The odd-shift corollary is a local-parity sanity check: the target-free
model does not encode the singular local factor.  For admissible even
shifts, replacing the explicit odd-shift value by
$(\mathfrak S(h)-1)XI_0(W)$ is exactly the unproved prime-pair step.  No
such replacement is made here.
\end{remark}
